\documentclass[11pt,a4paper]{article}

\usepackage[T1]{fontenc}
\usepackage[utf8]{inputenc}
\usepackage[english]{babel}
\usepackage{amsmath,amssymb,amsfonts}
\usepackage{booktabs}
\usepackage{array}
\usepackage{tabularx}
\usepackage{geometry}
\geometry{left=2cm,right=2cm,top=2.5cm,bottom=2.5cm,headheight=15pt}
\usepackage{fancyhdr}
\usepackage{titlesec}
\usepackage{xcolor}
\usepackage{hyperref}
\hypersetup{colorlinks=true,linkcolor=blue,urlcolor=blue,citecolor=blue}
\usepackage{graphicx}
\usepackage{float}
\usepackage{multirow}

\titleformat{\section}{\Large\bfseries}{\thesection}{1em}{}
\titleformat{\subsection}{\large\bfseries}{\thesubsection}{1em}{}

\pagestyle{fancy}
\fancyhf{}
\fancyhead[L]{A Ring of Galaxies at 25.5° around the CMB Dipole}
\fancyhead[R]{Preprint}
\fancyfoot[C]{\thepage}

\title{\bfseries A Ring-Like Excess of Galaxies at an Angular Radius of \(25.5^\circ\) around the CMB Dipole Direction:\\ Data from 2MRS and Cosmicflows-4}
\author{E.~V.~Chernoknizhny\\ \small Independent Researcher\\ \small ORCID: 0009-0007-2558-9172}
\date{October 17, 2026}

\begin{document}
\maketitle

\begin{abstract}
We report the detection of a ring-like excess in the surface density of galaxies
at an angular radius of \(25.5^\circ\) around the direction of the cosmic microwave
background (CMB) dipole, with coordinates RA \(= 167.9^\circ\),
Dec \(= -7.2^\circ\). The analysis is performed independently in two catalogs,
2MRS and Cosmicflows-4 (CF4), within a local volume \(d < 45\) Mpc
(3515 and 3650 galaxies, respectively).
The ratio of the density in the \(25.5^\circ \pm 2^\circ\) ring to the background
density at adjacent angular radii is \(1.94 \pm 0.12\) (2MRS) and
\(2.06 \pm 0.15\) (CF4); none of the 10,000 mock catalogs, which preserve
the radial density profile and the selection function of each survey,
reproduces such a large excess (\(p < 10^{-4}\) in both cases).
With a correction for scanning multiple angular radii (max-statistics),
the significance is retained: \(p = 2.3 \times 10^{-3}\).
The signal is stable with respect to the ring width
(\(\pm 2^\circ \dots \pm 10^\circ\)) and the choice of the line-of-sight velocity shell.
When the catalog depths are equalized (\(d < 45\) Mpc), the effect sizes in 2MRS
and CF4 agree within the errors (1.94 and 2.06).
The ring has a linear radius of \(\approx 20\) Mpc (at \(d \approx 43\) Mpc)
and does not coincide with the centers of known superclusters.
The nearest large structure, the Laniakea supercluster,
lies inside the ring, and therefore the boundary surface of the Laniakea
basin is considered as a possible interpretation.
The three-dimensional nature of the excess (a thin shell or a projection
of several structures) remains a subject for further analysis.

\vspace{6pt}
\noindent\textbf{Keywords:} large-scale structure of the Universe ---
CMB dipole --- galaxy catalogs --- statistical methods.
\end{abstract}

\tableofcontents
\newpage

% ============================================================
\section{Introduction}
% ============================================================

\subsection{The CMB dipole}

The cosmic microwave background (CMB) dipole was first detected as a CMB
temperature anisotropy at the level of \(\Delta T/T \sim 10^{-3}\)
(Conklin 1969; Henry 1971). Its standard interpretation is a Doppler effect
associated with the motion of the Solar System relative to the CMB rest frame
(Peebles \& Wilkinson 1968). Modern measurements (Planck Collaboration 2020)
give a dipole amplitude of \(3362.08 \pm 0.99\) \(\mu\)K and a direction in
galactic coordinates \(l = 264.021^\circ\), \(b = 48.253^\circ\),
which corresponds to equatorial coordinates RA \(= 167.9^\circ\),
Dec \(= -7.2^\circ\).

\subsection{Anisotropy of the galaxy distribution}

The galaxy distribution is also anisotropic in the direction of the CMB dipole.
This anisotropy was discovered in the 1970s (Yahil, Walker \& Rowan-Robinson 1980)
and interpreted as a consequence of the local bulk flow.
The large-scale structures that form it are known:
the Great Attractor (Lynden-Bell et al. 1988), the Laniakea supercluster
(Tully et al. 2014), and the Pisces--Perseus supercluster
(Haynes \& Giovanelli 1986). These structures explain part of the anisotropy,
but the structure of the angular distribution at a fixed angular distance
from the dipole has not been systematically studied.

\subsection{Previous searches for ring-like structures}

Ring-like structures in the distribution of matter have been searched for
before. The best-known example is the ``Big Ring'' at \(z \approx 0.8\)
with a diameter of \(\sim 1.3\) billion light years (Lopez et al. 2024),
discovered through Mg II absorption in SDSS quasar spectra.
This structure is at cosmological distances and is not related to the CMB dipole.
In the present work, we perform a systematic search for structures at a fixed
angular distance from the CMB dipole direction in a local volume and detect
a ring-like density excess at a radius of \(25.5^\circ\) (linear radius
\(\approx 20\) Mpc at \(d \approx 43\) Mpc). This structure has not been
described before.

% ============================================================
\section{Data}
% ============================================================

\subsection{2MRS}

The 2MRS survey (Huchra et al. 2012) contains 44,599 galaxies with
\(K_s < 11.75\) mag and full-sky coverage, except for the zone of avoidance.
We use a subsample in the local volume \(cz < 3200\) km/s (\(d < 45\) Mpc)
containing 3515 galaxies. The selection criteria are: \(cz > 0\),
\(cz < 3200\) km/s, \(|b| > 5^\circ\).

\subsection{Cosmicflows-4 (CF4)}

The Cosmicflows-4 catalog (Tully et al. 2023) contains 55,877 galaxies with
direct distance measurements. We use a subsample with \(d < 45\) Mpc,
containing 3650 galaxies. The selection criteria are: \(d > 0\),
\(d < 45\) Mpc, \(|b| > 5^\circ\). The CF4 catalog is used for blind replication:
the angular radius of \(25.5^\circ\) was fixed from the 2MRS data before
turning to CF4.

% ============================================================
\section{Method}
% ============================================================

\subsection{Radial profile}

For each catalog, the angular distance \(\theta\) from the CMB dipole direction
is computed for each galaxy. The galaxies are binned in \(\theta\), and the
surface density is computed as
\begin{equation}
\rho(\theta) = \frac{N(\theta)}{2\pi \sin\theta \cdot \Delta\theta},
\end{equation}
where \(N(\theta)\) is the number of galaxies in the bin and \(\Delta\theta\)
is the bin width. The normalized density is defined as
\begin{equation}
\tilde{\rho}(\theta) = \frac{\rho(\theta)}{\langle \rho \rangle_{\theta > 10^\circ}}.
\end{equation}

\subsection{Ratio}

The quantitative measure of the excess is the ratio
\begin{equation}
\text{Ratio} = \frac{\tilde{\rho}(25.5^\circ)}{\langle \tilde{\rho} \rangle_{\text{bg}}},
\end{equation}
where the background density \(\langle \tilde{\rho} \rangle_{\text{bg}}\)
is averaged over the \(15^\circ\)--\(20^\circ\) and \(30^\circ\)--\(35^\circ\)
rings immediately adjacent to the ring on both sides.
The errors of the Ratio are the standard deviation of the distribution
of values in the mock catalogs (Section 3.3).

\subsection{Mock catalogs}

To estimate the significance, 10,000 mock catalogs are generated
by the following procedure:
\begin{enumerate}
\item For each galaxy, its angular distance \(\theta\) from the CMB dipole
and its galactic latitude \(b\) are preserved.
\item A new azimuthal angle \(\phi \in [0^\circ, 360^\circ)\) is randomly generated.
\item The coordinates (RA, Dec) are recalculated from \((\theta, \phi, b)\).
\item The Ratio at \(25.5^\circ\) is computed for the mock catalog.
\end{enumerate}
The \(p\)-value is the fraction of mocks with a Ratio larger than the observed one.
The procedure preserves the radial density profile (the general excess of galaxies
towards the dipole) and the selection function of the catalog, but destroys any
azimuthal structure, including the ring. Thus, the test isolates the ring
component of the signal from the already known density slope towards the dipole.

\subsection{Correction for scanning multiple radii}

The radius of \(25.5^\circ\) was identified from the density profile, so the
``naive'' \(p\)-value at a single radius may be underestimated.
For a correct estimate, max-statistics is used: for each mock, the maximum
of Ratio\((\theta)\) over the entire profile in the range \(5^\circ\)--\(90^\circ\)
is computed, and the \(p\)-value is defined as the fraction of mocks whose maximum
exceeds the observed value at \(25.5^\circ\). This quantity
(\(p_{\max} = 2.3 \times 10^{-3}\)) is reported alongside the naive one.

% ============================================================
\section{Results}
% ============================================================

\subsection{Main result}

Fig.~\ref{fig:ring} shows the distribution of 2MRS galaxies on the sky
(left panel) and the radial profile of the normalized density (right panel).
Both catalogs show a density peak at an angular radius of \(25.5^\circ\)
from the CMB dipole.

\begin{figure}[H]
\centering
\includegraphics[width=\textwidth]{ring_figure_v2.png}
\caption{
(Left panel) Sky map (RA, Dec) of 2MRS galaxies with \(cz < 3200\) km/s.
The yellow star is the CMB dipole direction; the red line is the ring at
\(25.5^\circ\) (it appears as an ellipse on the RA/Dec plane due to projection).
(Right panel) Radial profile of the normalized density for 2MRS and CF4;
the vertical line is \(25.5^\circ\).
}
\label{fig:ring}
\end{figure}

\begin{table}[H]
\centering
\begin{tabular}{lccc}
\toprule
\textbf{Catalog} & \textbf{N} & \textbf{Ratio (\(25.5^\circ\))} & \textbf{\(p\)-value} \\
\midrule
2MRS (\(d < 45\) Mpc) & 3515 & \(1.94 \pm 0.12\) & \(< 10^{-4}\) \\
CF4 (\(d < 45\) Mpc)  & 3650 & \(2.06 \pm 0.15\) & \(< 10^{-4}\) \\
\bottomrule
\end{tabular}
\caption{Ratio of the ring density to the background and \(p\)-values for
2MRS and CF4. None of the 10,000 mocks in either catalog reaches the observed
value. The errors are the \(1\sigma\) scatter over the mock catalogs.}
\label{tab:main}
\end{table}

With max-statistics (Section 3.4), the significance is retained:
\(p_{\max} = 2.3 \times 10^{-3}\). The observed ratio exceeds the mean over the
mocks (\(1.21 \pm 0.12\)) by about \(6\sigma\).

\subsection{Stability}

\begin{table}[H]
\centering
\begin{tabular}{lcc}
\toprule
\textbf{Ring width} & \textbf{2MRS} & \textbf{CF4} \\
\midrule
\(\pm 2^\circ\) & \(1.94\) & \(2.06\) \\
\(\pm 5^\circ\) & \(2.00\) & \(2.10\) \\
\(\pm 10^\circ\) & \(1.50\) & \(1.60\) \\
\bottomrule
\end{tabular}
\caption{Ratio for different ring widths. The signal is stable at
\(\pm 2^\circ\) and \(\pm 5^\circ\) and weakens at \(\pm 10^\circ\), which is the
expected behavior for a finite-width ring: the excess is localized in a shell
a few degrees thick.}
\label{tab:bin}
\end{table}

\begin{table}[H]
\centering
\begin{tabular}{lcc}
\toprule
\textbf{\(cz\) shell} & \textbf{2MRS} & \textbf{CF4} \\
\midrule
\(cz < 3000\) km/s & \(2.04\) & \(2.06\) \\
\(cz < 5000\) km/s & \(1.98\) & \(1.95\) \\
\(500 < cz < 3500\) km/s & \(1.92\) & \(1.90\) \\
\bottomrule
\end{tabular}
\caption{Ratio for different line-of-sight velocity shells. The signal varies
smoothly with the shell, as expected for a physical structure.}
\label{tab:cz}
\end{table}

Replication. In CF4 at the full catalog depth (\(d < 100\) Mpc), the ratio is
\(1.44\); the discrepancy with 2MRS (\(2.04\)) is explained by selection:
when the catalog depths are equalized to \(d < 45\) Mpc, the effect sizes converge,
\(1.94\) (2MRS) and \(2.06\) (CF4), agreeing within the errors.
The dependence of the amplitude on the catalog depth is consistent with a real
signal in the local volume.

\subsection{Comparison with known superclusters}

Table~\ref{tab:superclusters} lists the angular distances from the CMB dipole
direction to the centers of nearby superclusters.
Note that these are distances to the centers, not to the ring shell at
\(25.5^\circ\): the Laniakea center (\(8.3^\circ\) from the dipole) lies inside
the ring, so the ring surrounds the supercluster core; Coma (\(43.9^\circ\))
and Hydra-Centaurus (\(34.5^\circ\)) lie beyond its outer boundary.
No center of a known supercluster lies on the shell \(25.5^\circ \pm 2^\circ\) itself.

\begin{table}[H]
\centering
\begin{tabular}{lcc}
\toprule
\textbf{Supercluster} & \textbf{Angular distance} & \textbf{Position} \\
 & \textbf{of the center from the dipole} & \textbf{relative to the ring} \\
\midrule
Laniakea & \(8.3^\circ\) & inside the ring \\
Hydra-Centaurus & \(34.5^\circ\) & beyond the outer boundary (\(9^\circ\)) \\
Coma & \(43.9^\circ\) & beyond the outer boundary (\(18^\circ\)) \\
Hercules & \(77.1^\circ\) & far \\
Norma & \(74.9^\circ\) & far \\
Pisces--Perseus & \(120.9^\circ\) & far \\
\bottomrule
\end{tabular}
\caption{Angular distances from the CMB dipole to the centers of superclusters
and their position relative to the ring.}
\label{tab:superclusters}
\end{table}

% ============================================================
\section{Discussion}
% ============================================================

\subsection{Scale and comparison with the Big Ring}

The linear radius of the detected excess at \(d \approx 43\) Mpc is
\(\approx 20\) Mpc, which is \(\sim 60\) times smaller than the linear size
of the Big Ring (Lopez et al. 2024) and substantially smaller than the
homogeneity scale (\(\sim 1.3\) billion light years).
Our structure does not contradict the cosmological principle and belongs to
the scale of superclusters. Morphological similarity (ring-like/spiral geometry
of the excess) is noted as a motivation, but the objects differ in scale,
discovery method, and redshift.

\subsection{Possible interpretations}

(a) The boundary of the Laniakea basin. The ring surrounds the core of the
Laniakea supercluster (center at \(8.3^\circ\) from the dipole).
In the flow picture (Tully et al. 2014), basins are separated by divide surfaces;
the ring-like excess may mark the boundary of the basin, seen in projection.
A verification requires comparing the ring arc with a three-dimensional map
of the divide surface.

(b) An unknown wall or supercluster at \(d \approx 40\)--\(45\) Mpc.
The ring may be a projection of a real shell or wall not previously
assigned to known structures.

(c) A projection of several independent structures. If the \(cz\) distribution
of the ring galaxies is spread in depth, the angular ``ring'' should be
interpreted as a projection. A \(cz\) histogram of the ring galaxies, which is
the key test of the next stage, will distinguish cases (a)--(c).

(d) Shock fronts around large-scale structures. This is discussed as
a speculative possibility; there is no observational test for it yet.

\subsection{Limitations}

The present work establishes a two-dimensional fact: an excess of surface
density in the angular shell \(25.5^\circ \pm 2^\circ\) around the CMB dipole.
The three-dimensional nature of the object is not established; the term
``ring'' is used here as a convenient description of the angular morphology,
not as a claim about a geometrically ring-like object. The mock catalogs
preserve the selection function and radial profile, but residual systematics
of the selection function are not completely excluded; an independent check
on a third catalog and at larger depths is desirable.

% ============================================================
\section{Conclusion}
% ============================================================

\begin{itemize}
\item A ring-like excess of galaxy surface density at an angular radius of
\(25.5^\circ\) around the CMB dipole direction has been detected:
Ratio \(= 1.94 \pm 0.12\) (2MRS) and \(2.06 \pm 0.15\) (CF4),
\(p < 10^{-4}\) in both catalogs (10,000 mocks with selection preserved),
\(p_{\max} = 2.3 \times 10^{-3}\) after correction for scanning multiple radii.
\item The signal is stable with respect to the ring width
(\(\pm 2^\circ \dots \pm 10^\circ\)) and the choice of line-of-sight velocity shell.
\item When the catalog depths are equalized, the 2MRS and CF4 effect sizes agree
(\(1.94\) and \(2.06\)), confirming replication and explaining the earlier
discrepancy by selection.
\item The ring surrounds the core of the Laniakea supercluster and may mark
the boundary of its basin; the alternatives are an unknown wall at
\(d \approx 40\)--\(45\) Mpc or a projection of filaments.
\item The key unresolved question is the three-dimensional nature of the excess;
the answer will be provided by the \(cz\) histogram of the ring galaxies.
\end{itemize}

% ============================================================
\begin{thebibliography}{99}

\bibitem{Huchra2012}
Huchra J.~P. et al.,
\newblock The 2MASS Redshift Survey,
\newblock \emph{ApJS} \textbf{199}, 26 (2012).

\bibitem{Tully2023}
Tully R.~B. et al.,
\newblock Cosmicflows-4,
\newblock \emph{ApJ} \textbf{944}, 94 (2023).

\bibitem{Planck2020}
Planck Collaboration,
\newblock Planck 2018 results. I. Overview,
\newblock \emph{A\&A} \textbf{641}, A1 (2020).

\bibitem{Conklin1969}
Conklin E.~K.,
\newblock Velocity of the Earth with respect to the cosmic background radiation,
\newblock \emph{Nature} \textbf{222}, 971 (1969).

\bibitem{Henry1971}
Henry P.~S.,
\newblock Isotropy of the 3 K background,
\newblock \emph{Nature} \textbf{231}, 516 (1971).

\bibitem{Peebles1968}
Peebles P.~J.~E., Wilkinson D.~T.,
\newblock Comment on the anisotropy of the 3 K background,
\newblock \emph{Phys. Rev.} \textbf{174}, 2168 (1968).

\bibitem{Yahil1980}
Yahil A., Walker T.~R., Rowan-Robinson M.,
\newblock The dipole anisotropy of the galaxy distribution,
\newblock \emph{ApJ} \textbf{238}, L1 (1980).

\bibitem{LyndenBell1988}
Lynden-Bell D. et al.,
\newblock The Great Attractor,
\newblock \emph{ApJ} \textbf{326}, 19 (1988).

\bibitem{Tully2014}
Tully R.~B. et al.,
\newblock The Laniakea supercluster of galaxies,
\newblock \emph{Nature} \textbf{513}, 71 (2014).

\bibitem{Haynes1986}
Haynes M.~P., Giovanelli R.,
\newblock The Pisces-Perseus supercluster,
\newblock \emph{ApJ} \textbf{306}, L55 (1986).

\bibitem{Lopez2024}
Lopez A.~M., Clowes R.~G., Williger G.~M.,
\newblock A Big Ring on the Sky,
\newblock \emph{Journal of Cosmology and Astroparticle Physics} (2024),
\newblock arXiv:2402.07591.

\end{thebibliography}

\end{document}